\begin{table}[ht]
\centering
\caption{Fiscal Multipliers, Baseline Model}
\label{tab:mult_main}
\begin{tabular}{lcccc}
\hline\hline
Instrument & $H=1$ & $H=4$ & $H=12$ & $H=20$ \\
\hline
\multicolumn{5}{l}{\textit{Panel A. Baseline model (fiscal rules with debt and output feedback)}} \\
Government spending      & 1.12 & 1.04 & 0.42 & $-$0.26 \\
                         & [1.05, 1.18] & [0.90, 1.15] & [0.20, 0.65] & [$-$0.59, 0.03] \\
Labor tax                & 0.32 & 0.62 & 0.65 & 0.18 \\
                         & [0.26, 0.39] & [0.49, 0.76] & [0.45, 0.89] & [$-$0.10, 0.45] \\
Capital tax              & 1.35 & 2.13 & 2.82 & 2.70 \\
                         & [1.23, 1.47] & [1.90, 2.36] & [2.42, 3.27] & [2.23, 3.24] \\
Consumption tax          & 0.27 & 0.39 & 0.03 & $-$0.62 \\
                         & [0.21, 0.32] & [0.29, 0.48] & [$-$0.16, 0.25] & [$-$0.93, $-$0.31] \\
Investment tax credit    & 0.76 & 1.54 & 2.23 & 2.14 \\
                         & [0.59, 0.92] & [1.23, 1.86] & [1.71, 2.81] & [1.54, 2.86] \\
Expensing rate           & 0.20 & 0.40 & 0.19 & $-$0.37 \\
                         & [0.12, 0.29] & [0.23, 0.57] & [$-$0.02, 0.47] & [$-$0.71, $-$0.06] \\
\hline
\multicolumn{5}{l}{\textit{Panel B. \cn{Zubairy2014}, Table 2}} \\
Government spending      & 1.07 & 1.06 & 0.90 & 0.72 \\
                         & [1.01, 1.13] & [0.96, 1.17] & [0.73, 1.07] & [0.49, 0.93] \\
Labor tax                & 0.13 & 0.32 & 0.68 & 0.85 \\
                         & [0.09, 0.18] & [0.22, 0.45] & [0.41, 1.09] & [0.42, 1.58] \\
Capital tax              & 0.34 & 0.43 & 0.52 & 0.46 \\
                         & [0.30, 0.37] & [0.34, 0.51] & [0.30, 0.73] & [0.15, 0.81] \\
\hline\hline
\multicolumn{5}{p{0.85\textwidth}}{\footnotesize Note: Present-value own-revenue multipliers (equation~\ref{eq:mult_main}). Posterior median with 5--95 percentile intervals in brackets, from 1{,}000 posterior draws. A positive value means output moves in the expansionary direction of the policy. Sample 1961Q1--2013Q4. Panel B reproduces \cn{Zubairy2014}; her tax multipliers use the change in total tax revenue as the denominator, so they are not directly comparable to Panel A.}
\end{tabular}
\end{table}
